% =====================================================================
% Additional experiment: MRL beyond English
% Code and results: experiments/indic_experiment/
% =====================================================================
\section{Additional Experiment: Does MRL Transfer Beyond English?}
\label{sec:additional-indic}

\paragraph{Motivation and research question.}
The five main experiments evaluate MIPIC and its components primarily on
English-language data.  This leaves open whether the benefit of Matryoshka
Representation Learning (MRL) is tied to English-specific linguistic
structure, or whether it also transfers to Indic languages. The question is whether MRL improves
low-dimensional sentence representations for Indic sentiment classification
relative to the same frozen TinyBERT encoder.

\paragraph{Data.}
We use Hindi (\texttt{hi}), Bengali (\texttt{bn}), Gujarati (\texttt{gu}),
Tamil (\texttt{ta}), and Telugu (\texttt{te}).  The unsupervised MRL corpus
contains 20,240 Indic sentences.  For each language, the downstream probe
uses 156 labeled training examples and 998 held-out test examples.  The
unlabeled MRL corpus is kept separate from the labeled sentiment splits.
Thus, the result measures transfer from representation training rather than
supervised training on the evaluation labels.

\paragraph{Models and training.}
The encoder is TinyBERT General 6L with a 768-dimensional hidden size
\citep{jiao-etal-2020-tinybert}.  We compare a frozen pretrained TinyBERT
against TinyBERT fine-tuned with the MRL objective.  MRL uses two independent
dropout views of the same sentence as a positive pair and a symmetric
InfoNCE loss at the nested dimensions
\[
\{16,32,64,128,256,512,768\}.
\]
The run uses five epochs, AdamW with learning rate $2\times10^{-5}$ and
weight decay $0.01$, batch size 32, maximum sequence length 128, and
temperature $0.05$.  The reported MRL representation is the final saved
checkpoint.  Sentence vectors are mean-pooled token representations.

\paragraph{Probe and metrics.}
For each language and dimension, a logistic-regression classifier is fitted
on the labeled training split and evaluated on the held-out test split.
Accuracy measures the fraction of correct predictions, while macro-F1 gives
equal weight to the three sentiment classes and is therefore more informative
when class frequencies are imbalanced.  The main comparison is MRL minus the
frozen baseline at the same dimension and language.

\paragraph{Decision rule.}
Because this is an exploratory additional experiment and uses one training
run, we do not claim a statistical significance threshold or a preregistered
pass/fail result.  We regard the cross-lingual result as positive if MRL
outperforms frozen TinyBERT in macro-F1 at the full 768 dimensions for all
five languages, and if its mean macro-F1 is higher at every tested dimension.

\begin{table}[t]
\centering
\small
\setlength{\tabcolsep}{3.5pt}
\begin{tabular}{@{}lrrrrrrr@{}}
\toprule
Dimension & 16 & 32 & 64 & 128 & 256 & 512 & 768 \\
\midrule
Frozen accuracy & .493 & .493 & .493 & .493 & .493 & .504 & .531 \\
MRL accuracy & .527 & .536 & .539 & .549 & .562 & .571 & .572 \\
Frozen macro-F1 & .330 & .330 & .330 & .330 & .331 & .362 & .427 \\
MRL macro-F1 & .455 & .489 & .508 & .530 & .552 & .565 & .569 \\
\bottomrule
\end{tabular}
\caption{Mean IndicSentiment performance over Hindi, Bengali, Gujarati,
Tamil, and Telugu. Values are averaged over the five language-specific test
sets; the probe is refit separately for each language and dimension.}
\label{tab:indic-mrl}
\end{table}

\paragraph{Results.}
MRL improves the mean macro-F1 at every dimension, from $0.330$ for frozen
TinyBERT to $0.455$ at 16 dimensions and $0.569$ at 768 dimensions
(Table~\ref{tab:indic-mrl}).  Mean accuracy similarly increases from
$0.493$ for the frozen model at 16--128 dimensions to $0.527$--$0.549$ for
MRL, and reaches $0.572$ at 768 dimensions.  The improvement is not limited
to the largest representation: the average macro-F1 gain is approximately
12.5 points at 16 dimensions and remains approximately 14.2 points at 768
dimensions.

At 768 dimensions, the MRL macro-F1 gains over frozen TinyBERT are
22.3 points for Hindi, 13.4 for Bengali, 14.7 for Gujarati, 14.0 for Tamil,
and 6.5 for Telugu.  Thus, the direction of the effect is consistent across
all five languages, although its size varies substantially.  The strongest
full-width result is Hindi, while Telugu shows the smallest gain.

\paragraph{Takeaway.}
This additional experiment provides preliminary evidence that MRL is not
effective only because of English-specific hierarchies: it improves Indic
sentiment probes across five languages and at compact as well as full
dimensions.  However, the experiment has only one training seed, small
language-specific labeled splits, and a frozen TinyBERT baseline that is not
Indic-specialized.  The result should therefore be treated as evidence of
transferability, not as a controlled claim that MRL is language-independent.
